\section{Testbenches}
Hieronder zijn de testbenches te vinden. De testbenches zijn gegenereerd met AI. 

\subsection{Buzzer}
De testbench voert hier de volgende tests uit: \newline
1. Er wordt gekeken of er de buzzer aan gaat zonder invoer. Dit slaagt als de buzzer niet aangaat. \newline
2 t/m 5. De bestemde functies worden getest: 1 korte piep voor invoer, twee korte piepjes voor een correct wachtwoord, en een lange piep als de code fout is, of als de reset-knop wordt ingedrukt. \newline
6. Er wordt getest of de buzzer langer aangaat dan zou moeten wanneer een knop lang ingedrukt blijft. Test slaagt als de buzzer uitgaat na de aangegeven lengte. \newline
7. De prioriteit van de pulses wordt getest door alle pulsen hoog te maken. De prioriteit moet correct > fout > reset > press zijn. Test slaagt als de buzzer ,bij meerdere pulses, de puls met de hoogste prioriteit volgt. \newline
\lstinputlisting[caption=buzzer\_tb, style=vhdl]{../code/tb/buzzer\_tb.vhd}

\subsection{Clock Divider}
Er wordt gekeken of de clockOut-variabele hoog wordt na elke vier clock cycles. Dit is normaal elke 25 miljoen, maar om het testproces te versimpelen is een lager getal gekozen. \newline
\lstinputlisting[caption=ClockDivider\_tb, style=vhdl]{../code/tb/ClockDivider\_tb.vhd}

\subsection{Keypad}
Hier wordt getest of de invoer vanuit de keypad goed is vertaald naar een integer. Ook wordt de resetknop "*" getest. \newline
\lstinputlisting[caption=keypad\_tb, style=vhdl]{../code/tb/keypad\_tb.vhd}

\subsection{Passcode}
De testbench voert hier de volgende tests uit: \newline
1. Er wordt gekeken of bij een correct wachtword het signaal wordt doorgegeven dat het wachtwoord klopt. \newline
2. Er wordt gekeken of het systeem goed wordt gereset. \newline
3. Er wordt gekeken of bij een fout wachtword het signaal wordt doorgegeven dat het wachtwoord niet klopt. \newline
4. Dit is een combinatie van test 1 \& 2, waarbij het systeem eerst wordt gereset, waarna het juiste wachtwoord wordt ingevoerd. \newline
\lstinputlisting[caption=passcode\_tb, style=vhdl]{../code/tb/passcode\_tb.vhd}

\subsection{Seven Segment Driver}
De testbench checkt eerst of de displays uit staan na een reset-signaal. Hierna wordt er getest of er bij elke nieuwe invoer een nieuw streepje tevoorschijn komt bij de displays. Dit begint bij HEX0 en de laatste zou HEX3 moeten zijn. \newline
\lstinputlisting[caption=seven\_segment\_driver\_tb, style=vhdl]{../code/tb/seven\_segment\_driver\_tb.vhd}

\subsection{State Machine}
De testbench voert hier, voor elke transitie naar een andere staat, een test uit om te bepalen of deze daadwerkelijk gebeurt bij de bijbehorende waarden. \newline
\lstinputlisting[caption=statemachine\_tb, style=vhdl]{../code/tb/statemachine\_tb.vhd}

